\section{Conclusion}
\label{sec:conclusion}

AutoHeal integrates Papers 3--9 of the VulnPrio sequence into a
seven-stage self-healing pipeline with pre-registered safety bounds.
The 2{,}700-row frozen evaluation reports honestly on which regimes
admit autonomous remediation and which do not. The substantive
findings appear in \S\ref{sec:results}; the pre-registered
hypothesis outcomes appear in Table~\ref{tab:hypotheses}.

\textbf{Synthesis claim.} The components developed in
Papers 3--9 are sufficient to build a self-healing framework whose
\emph{safety bounds are observable in advance} and whose
\emph{failure modes are predicted by the program's theoretical
results} (Paper~9 Theorem~1 predicts the K=200 exclusion;
Paper~7's lag-1 hazard predicts the calibration recipe choice).
Building AutoHeal does not require new theoretical machinery beyond
what the prior papers established; it requires \emph{integrating}
them under a pre-registered architecture.

\textbf{Honest scope.} AutoHeal is evaluated on the EEHDA synthetic
fleet with real CVE attribute distributions. The pre-registered
failure-mode distribution (92/5/3) is drawn from public literature
rather than measured on a real fleet. The cell-mean numbers reported
here may not transfer; the qualitative findings (which capacity
regimes admit autonomous remediation, where safety bounds fire) are
designed to be transferable to a real-deployment evaluation
following the same pre-registration discipline.

\textbf{What this paper adds to the program.} AutoHeal is the
\emph{integration} paper of the VulnPrio sequence. Papers 1--9
established components and bounds; AutoHeal demonstrates that the
components fit together in an architecture whose behaviour is
predicted by the prior papers' findings. The program's
recommendations now have a concrete object to attach to: ``deploy
AutoHeal at $K \leq 100$ with the pre-registered triage thresholds
and safety bounds; do not deploy at $K = 200$ per Paper~9
Theorem~1 and Corollary~4.''

\textbf{Reproducibility.} The framework, runner, analysis, and
frozen results are available at
\url{https://github.com/Harshavardhanmalla-Labs/research-papers/tree/main/paper10}.

\textbf{Future work.} The most consequential remaining direction is
external validation on real fleet telemetry. AutoHeal's
pre-registered architecture and safety bounds are designed to be
directly portable to a real-deployment study; the failure-mode
distribution would be measured from the deployment rather than
pre-registered, and the cell-mean numbers would be re-estimated.
Closed-loop calibration of the triage thresholds (treating the
thresholds themselves as parameters to learn rather than constants
to pre-register) is the second most consequential direction, but
requires careful pre-registration discipline to avoid
post-hoc tuning.
